\documentclass{beamer}
\usetheme{default}
\beamertemplatenavigationsymbolsempty
\usepackage{rembeamer}

\begin{document}

\begin{frame}
  \begin{itemize}
    \item[] 
      Explaining how we can use randomization for hypothesis testing
      and interpret the result, including the p-value.
  \end{itemize}
\end{frame}

\begin{frame}
  \frametitle{Frequency Table: Gender and Decision}
  \begin{table}
    \begin{tabular}{@{}rlll@{}}
      \toprule
      gender & promoted & not promoted & total \\ 
      \cmidrule(r){1-1} \cmidrule(lr){2-2} \cmidrule(lr){3-3}
      \cmidrule(l){4-4}
      men & 21 & 3 & 24 \\[0.3em]
      women & 14 & 10 & 24 \\[0.3em]
      \cmidrule(r){1-1} \cmidrule(lr){2-2} \cmidrule(lr){3-3}
      \cmidrule(l){4-4}
      total & 35 & 13 & 48 \\[0.3em]
      \bottomrule
    \end{tabular}
  \end{table}
  \begin{itemize}
    \item[] Propotion of men promoted: $\dfrac{21}{24} = 0.875$. 
      \vspace{0.5cm}
    \item[] Propotion of women promoted: $\dfrac{14}{24} = 0.583$. 
      \vspace{0.5cm}
    \item[] Difference in promotion proportions: $0.875 - 0.583 = 0.292$.
  \end{itemize}
\end{frame}

\begin{frame}
  \frametitle{Hypotheses}
  \begin{itemize}
    \item[] $H_0$: Null hypothesis. The variables `gender' and
    `decision' are independent. 
      \vspace{0.5cm}
    \item[] $H_0$: $p_m - p_f = 0$. 
      \vspace{0.5cm}
    \item[] $H_A$: Alternative hypothesis. The variables `gender' and
    `decision' are not independent. 
      \vspace{0.5cm}
    \item[] $H_A$: $p_m - p_f > 0$
      \vspace{0.5cm}
    \item[] $H_A$: $p_m - p_f \neq 0$
  \end{itemize}
\end{frame}

\begin{frame}
  \frametitle{Randomization}
  \begin{itemize}
    \item[] Total promotion rate: 35 promoted, 13 not promoted. 
      \vspace{0.5cm}
    \item[] Applicants are 24 men and 24 women.
      \vspace{0.5cm}
    \item[] Randomly assign promotions to the applicants: 35 randomly
    split up between 24 men and 24 women.
  \end{itemize}
\end{frame}

\begin{frame}
  \frametitle{Random Trials}
  \begin{displaymath}
    \begin{tabular}{@{}ccc@{}}
      \toprule
      men promoted & women promoted & difference \\[0.3em]
      \cmidrule(r){1-1} \cmidrule(lr){2-2} \cmidrule(lr){3-3}
      20 & 15 & 0.208 \\
      22 & 13 & 0.357 \\
      15 & 20 & -0.208 \\ 
      20 & 15 & 0.208 \\
      11 & 24 & -0.542 \\  
      19 & 16 & 0.125 \\
      20 & 15 & 0.208 \\ 
      14 & 21 & -0.292 \\
      18 & 17 & 0.042 \\  
      12 & 23 & -0.458 \\
      \bottomrule
    \end{tabular}
  \end{displaymath} 
\end{frame}

\begin{frame}
  \frametitle{p-value}
  \begin{itemize}
    \item[] 
      Many, many random trials. 
      \vspace{0.5cm}
    \item[] 
      Promotion proportion difference $\geq 0.292$ happened about
      $2\%$ of the time. $p = 0.02$. 
      \vspace{0.5cm}
    \item[] 
      $2\%$ change that $H_0$ is true and our sample statistics was
      just random variance in the population. 
      \vspace{0.5cm}
    \item[] 
      Significance level $\alpha$. 
      \vspace{0.5cm}
    \item[] 
      $p < \alpha$ is called `statistically significant'. 
      \vspace{0.5cm}
    \item[] 
      Often $\alpha = 0.05$. If so, our sample statistics is
      significant to reject $H_0$.
    \item[] 
  \end{itemize}
\end{frame}

\end{document}


